\section{Resumen del trabajo realizado}
El trabajo expuesto a lo largo de este documento resulta de la consecución de cuatro fases concretas.

Conceptualmente se inicia con la contextualización del problema, donde se introduce la idea de buscar la forma de simular
la operación de una línea de autobús a lo largo de un día de servicio. Para ello, previa investigación de literatura existente,
se escoge el paradigma multi-agente, debido a las características del entorno dinámicas y complejas, junto a la capacidad de evolución
que presenta.

Una vez queda definido el objetivo, se empieza a trabajar sobre el conjunto de datos otorgado por la EMT. Durante esta fase, se acota
el alcance del simulador debido a distintas características observadas, como la exclusión de fines de semana o de días con gratuidad.
Debido también a la estructura de los datos, se caracteriza estadísticamente mediante distribuciones tanto valores comparativos 
de ocupación, como la carga y descarga de pasajeros, siendo estos últimos los principales parámetros estocásticos del sistema multi-agente.

Con los datos estudiados y el objetivo contextualizado, se inicia la fase de construcción del sistema. En esta fase se opta por el
desarrollo de dos sistemas paralelos, que se diferencian por la caracterización en uno de ellos de los pasajeros como agentes, pero
con el objetivo de que ambos sean plenamente funcionales bajo el escenario a modelar. En esta fase se modela la arquitectura
en la plataforma de diseño de agentes GAMA, y se implementan conceptos como la reconstrucción lineal conservativa, para incorporar
la tendencia en la generación de pasajeros, se diseñan protocolos de comunicación basados en los estándares FIPA, 
o se resuelven conflictos como la abstracción de la inferencia de bajada de pasajeros mediante acciones del entorno, extrayendo
estas capacidades de los agentes.

Por último, y en concordancia con lo establecido en los requisitos del sistema, se construye
una interfaz que facilita el lanzamiento y la visualización de las simulaciones, permitiendo
abordar la validación de forma más ágil.

Dicha validación muestra que el simulador reproduce la operativa de la línea con un sesgo
medio en torno a los dos pasajeros y una sobredispersión que crece a lo largo del recorrido,
magnitudes ambas caracterizadas y documentadas. Se identifica asimismo que el modelado de
descensos como eventos independientes restringe la validez del sistema al escenario base,
excluyendo el análisis de configuraciones alternativas de servicio.

\section{Satisfacción de objetivos}
Para el presente trabajo se establece un objetivo principal y tres objetivos secundarios. A lo largo del documento se ha discutido,
especialmente en el capítulo correspondiente a la validación y resultados (véase capítulo~\ref{chapter:application_cases}), la completitud
de los distintos objetivos. A continuación, para cada uno se formaliza y se discute en profundidad su grado de conformidad alcanzado.

\subsection{Satisfacción del objetivo principal del trabajo}
Este objetivo concretamente es \textbf{determinar en qué medida un sistema multi-agente parametrizado a partir de datos reales
es capaz de reproducir la operativa de una línea de autobús urbano.} Tras la construcción
y evaluación del sistema multi-agente para las dos implementaciones, se puede concluir que el objetivo principal del trabajo
resulta satisfecho, ya que se constata que, pese a la discusión relativa al sesgo y a la sobredispersión, el sistema es capaz
de capturar el comportamiento general de la línea de estudio para los días laborables y estudiar, en base a los indicadores
de rendimiento obtenidos, la capacidad de absorción de la demanda por parte de la operativa de servicio promedio.

\subsection{Satisfacción de los objetivos secundarios}

En cuanto a los objetivos secundarios, el primero de ellos, \textbf{caracterizar estadísticamente los datos operacionales de la línea},
se satisface mediante el análisis
recogido en el capítulo~\ref{chapter:param_analisys}. Los datos operacionales cedidos por la
EMT han sido depurados y caracterizados estadísticamente, ajustando distribuciones de
probabilidad sobre los valores de ocupación, embarque y descenso para cada combinación de
parada y franja horaria, con la selección de familia realizada mediante criterios de bondad
de ajuste y de información.

El segundo objetivo, \textbf{validar que ambas propuestas arquitectónicas generan resultados equivalentes},
queda igualmente cumplido. El contraste directo de los indicadores de
rendimiento generados por ambas soluciones, bajo idénticos parámetros de entrada, confirma que las dos propuestas arquitectónicas
producen resultados equivalentes. Se puede concluir entonces que tanto la representación agregada como la
representación basada en agentes pasajero individuales reproducen el comportamiento esperado.

Por último, el tercer objetivo, \textbf{proporcionar una herramienta que permita ejecutar, almacenar y visualizar de forma
interactiva las distintas simulaciones generadas por el sistema}, queda satisfecho mediante la construcción de la interfaz presentada
al inicio del capítulo~\ref{chapter:application_cases}, que unifica todos los requerimientos del objetivo.

\section{Conclusiones generales}
En un cómputo general, este trabajo de investigación se posiciona como la fase inicial o de validación de un proyecto colaborativo
entre entidades públicas, cuyo objetivo último es servir como herramienta para analizar y mejorar el servicio de autobuses 
urbanos.

El trabajo realizado asienta la base sobre la que evolucionar el sistema hacia dicho
objetivo. La elección del paradigma multi-agente permite un escalado posterior hacia
arquitecturas deliberativas sin alterar el esqueleto del modelo. La interfaz implementada
facilita el lanzamiento, almacenamiento y visualización de escenarios. Y el análisis de los
datos operacionales, si bien constituye el origen de la principal limitación identificada,
delimita con precisión qué aspectos del servicio resultan reproducibles y con qué grado de
fidelidad.

Cabe destacar que la validación del sistema no se ha limitado a demostrar el acuerdo entre
lo simulado y lo observado, sino que ha perseguido caracterizar dónde falla el modelo y por
qué mecanismo. Es dicha caracterización, más que el acuerdo en sí, la que permite establecer
sobre qué preguntas puede el simulador aportar respuesta y cuáles quedan condicionadas a
desarrollos posteriores.